% Part: lambda-calculus
% Chapter: lambda-definability
% Section: introduction

\documentclass[../../../include/open-logic-section]{subfiles}

\begin{document}

\olfileid[fr]{lam}{rep}{int}
\olsection{Introduction}

À première vue, le lambda-calcul n'est qu'un
calcul très abstrait d'expressions qui
représentent des fonctions et leurs applications
à d'autres fonctions. Sa syntaxe ne suggère
aucune espèce particulière d'objets auxquels
ces fonctions s'appliqueraient; elle ne
mentionne notamment pas les nombres naturels.
Ses opérations fondamentales, l'application
et l'abstraction lambda, concernent toute
fonction, pas seulement celles qui portent
sur les nombres naturels.

Avec un peu d'ingéniosité, on peut néanmoins
définir dans le lambda-calcul des fonctions
arithmétiques, c'est-à-dire des fonctions
sur les nombres naturels. Pour cela, on
associe à chaque entier naturel
$n\in\Nat$ un lambda-terme particulier
$\num{n}$, le \emph{numéral de Church}
de~$n$, nommé d'après Alonzo Church.

\begin{defn}
Pour $n\in\Nat$, le \emph{numéral de Church}
$\num{n}$ représente~$n$ :
\[
  \num{n} \ident \lambd[fx][f^n(x)].
\]
Ici, $f^n(x)$ désigne le résultat de
l'application de~$f$ à~$x$, $n$ fois.
Par exemple, $\num{0}$ est
$\lambd[fx][x]$ et $\num{3}$ est
$\lambd[fx][f(f(f\,x))]$.
\end{defn}

Le numéral $\num{n}$ représente ainsi
une fonction de deux arguments, $f$ et
$x$, qui renvoie $f^n(x)$. Les numéraux
de Church sont manifestement en forme
normale.

Cette représentation des nombres naturels
n'est utile que si elle permet de calculer.
Calculer avec les numéraux de Church
consiste à appliquer un lambda-terme~$F$
à un numéral~$\num{n}$, puis à réduire
$F\,\num{n}$ jusqu'à une forme normale.
Si, pour chaque~$n$, cette forme normale
existe et est un numéral~$\num{m}$,
on peut considérer~$m$ comme le résultat
du calcul et~$F$ comme la définition
d'une fonction $f\colon\Nat\to\Nat$ :
$f(n)=m$ si et seulement si
$F\,\num{n}\red\num{m}$. Par la
propriété de Church--Rosser, une forme
normale, si elle existe, est unique.
Ainsi, lorsque $F\,\num{n}\red\num{m}$,
ce terme ne peut se réduire à aucune
autre forme normale, notamment à aucun
autre numéral de Church.

Inversement, pour une fonction
$f\colon\Nat\to\Nat$ donnée, on peut
demander s'il existe un terme~$F$ qui
la définit ainsi. On dit alors que~$F$
\emph{$\lambda$-définit}~$f$ et que~$f$
est \emph{$\lambda$-définissable}.
Cette notion s'étend aux fonctions à
plusieurs arguments et aux fonctions
partielles.

\begin{defn}
Soit $f$ une fonction partielle de
$\Nat^k$ dans~$\Nat$. Un lambda-terme~$F$
\emph{$\lambda$-définit}~$f$ si, pour
tous $n_0$, \dots,~$n_{k-1}$,
\[
  F\,\num{n_0}\,\num{n_1}\dots\num{n_{k-1}}
  \red \num{f(n_0,n_1,\dots,n_{k-1})}
\]
lorsque $f(n_0,\dots,n_{k-1})$ est définie;
si cette valeur est indéfinie, le terme
$F\,\num{n_0}\,\num{n_1}\dots\num{n_{k-1}}$
n'a pas de forme normale.
\end{defn}

Les fonctions constantes fournissent
un exemple simple. Le terme
$C_k\ident\lambd[x][\num{k}]$
$\lambda$-définit la fonction
$c_k\colon\Nat\to\Nat$ telle que
$c_k(n)=k$. En effet, pour tout~$n$,
\[
  C_k\,\num{n}
  \ident(\lambd[x][\num{k}])\num{n}
  \redone\num{k}.
\]
La fonction identité est
$\lambda$-définie par $\lambd[x][x]$.
Définir des fonctions plus complexes
demande souvent davantage d'ingéniosité.
Il peut donc sembler surprenant que
toute fonction calculable soit
$\lambda$-définissable. La réciproque
vaut aussi : toute fonction
$\lambda$-définissable est calculable.

Pour une fonction partielle, la condition sur les entrées
sans valeur porte sur toutes les suites de réduction. Une
suite infinie peut partir d'un terme qui possède aussi une
forme normale. Une entrée sans valeur doit donner un terme
dont aucune suite de réductions n'atteint une forme normale.

\end{document}
